\section{模型假设与符号说明}

\subsection{模型假设}

假设 1：化学测量是每款啤酒的稳定属性，不随品评日期或品评场次变化。
化学量由仪器测定，与品评场次无关，每款啤酒因此对应唯一的化学谱。
据此化学谱可以按风格分层划分验证折，感官标签则须按日期组划分。
问题一以 250 款化学谱预测感官评分，问题二在化学与感官之间建立关联，问题三的生成模型按风格分层。

<<<<<<< HEAD
\begin{table}[htbp]
\centering\small
\caption{全局假设}
\label{tab:assumptions}
\begin{tabular}{>{\raggedright\arraybackslash}p{3.4cm}>{\raggedright\arraybackslash}p{4.6cm}>{\raggedright\arraybackslash}p{4.6cm}}
\toprule
假设 & 理由 & 对模型的影响 \\
\midrule
样本对齐以行号为准 & 原始评分文件中的编号对应附件一、四的行号，按行号对齐时逐维 Spearman 中位数为 0.969，按附件列对齐为 $-0.010$ & 所有连接都按行号，交叉验证分组与感官目标随之对齐 \\
\addlinespace
未校准信号列按 $\log_{10}$ 相对峰面积解释 & 157 个色谱质谱列没有绝对浓度校准数据，无法逐列证实底数 & 梯度提升树模型直接使用发布值，其分裂顺序不受底数选择影响；需要浓度时用 $10^x$ 还原 \\
\addlinespace
共同下限按未检出状态处理 & 142 列共 6,807 个单元等于共同下限 $-1.81781854$，现有材料不能确认仪器检出限 & 树模型把该值作为最低秩，线性模型加逐列下限指示变量 \\
\addlinespace
化学缺失值以训练折列均值填充 & 缺失只有 73 个单元，删去含缺失啤酒会损失 25.2\% 的样本 & 填充值只作计算占位，不解释为真实零或未检出 \\
\addlinespace
感官目标为当次出席品评员的平均分 & 加性模型中出席名单带来的偏移占啤酒间方差的中位数只有 2.2\% & 目标不校正出席名单差异，敏感性分析另行检验 \\
\addlinespace
细分风格标签准确 & 赛题以该标签作为问题三的条件变量 & 生成器按标签学习风格条件分布，不合并相近风格 \\
\addlinespace
文献气味阈值在啤酒基质中成立 & 阈值取自附件八的感官阈值 & 只有气味活度值的计算与迁移阈值依赖该假设 \\
\bottomrule
\end{tabular}
\end{table}
=======
假设 2：感官评分在同一品评日期组内相关，日期组是相关与独立性的单位。
同一日期组的啤酒共享品评员与品评情境，有各自的噪声、偏差与评分习惯。
因此交叉验证、重抽样与留出验证都以日期组为单位，问题一使用按日期组分组的嵌套交叉验证，问题二使用日期组簇 bootstrap，问题三的画像验证按日期组划分。
>>>>>>> 28f0f3593e345d7e82eeec8b19294d46c1357145

假设 3：分析范围限于观测到的化学取值，不向外推。
观测范围之外的取值没有数据支持，基于树的模型在范围外只能做常数外推。
因此生成与迁移输出的化合物谱被限制在观测范围内，超出范围的取值不作为结论依据。

假设 4：附件八的缺陷阈值单位与附件一对应列一致。
对数列的阈值取常用对数值，线性列的阈值取原值，阈值与化学列必须同尺度才可比。
这一假设决定化学层缺陷上界与超阈比例，迁移中另以不依赖阈值与单位的严格不增加规则作为对照。

\subsection{符号定义}

跨问题共用的符号列于表~\ref{tab:symbols}。
方法局部的符号在引入它们的一节定义，不在此列出。

\begin{table}[htbp]
\centering\small
\caption{共用符号}
\label{tab:symbols}
\begin{tabular}{cll}
\toprule
符号 & 含义 & 单位 \\
\midrule
$b$ & 啤酒索引，$b=1,\dots,250$ & 无 \\
$d$ & 感官维度，附件四的 50 列之一，$d=1,\dots,50$ & 无 \\
$X_b$ & 啤酒 $b$ 的化学特征，即附件一的 231 列 & 混合，见数据审计 \\
$\bar x_{bd}$ & 啤酒 $b$ 在维度 $d$ 上的品评小组原始评分均值 & 0--6 分 \\
$z_{bd}$ & 附件四标准化后的品评小组评分 & 以每位品评员标准差为单位 \\
$s(b)$ & 啤酒 $b$ 的细分风格标签，共 22 类 & 无 \\
\bottomrule
\end{tabular}
\end{table}

$\bar x_{bd}$ 是问题一与问题二的感官目标，也是问题三感官画像的观测真值。
$z_{bd}$ 是 $\bar x_{bd}$ 的确定性标准化，按品评员与维度标准化后再对小组取平均，去掉 0 到 6 量表的均值与标准差，只在数据核对、仅用附件的对照与问题二的敏感性分析中使用。
$s(b)$ 是问题三的条件变量，也是问题一基线所要超越的信息来源。
